
\needspace{2\baselineskip}
\color{uniblue}{\subsection{Контроль наличия напряжения (КНН)}}
\color{black}

\begin{enumerate}[label=\arabic{section}.\arabic{subsection}.\arabic*, labelsep=4pt, leftmargin=0pt, itemindent=57pt]

\item
КНН контролирует наличие напряжения на электрооборудовании. В устройстве контроль наличия напряжения представлен функцией «КНН».


\item
Функция содержит в себе ИО максимального напряжения, реагирующий на превышение минимального значения из подведенных линейных напряжений заданной уставки и ИО напряжения ОП, реагирующий на превышение напряжения ОП выше заданной уставки. Срабатывании ИО напряжения ОП блокирует пуск функции.

\item
Ввод функции в работу на этапе параметрирования устройства осуществляется программным переключателем SGF1 «Ввод\_функции». Информация о введенной функции отображается выходным сигналом «КНН: Ввод».

\item
Функция может быть выведена из работы оперативно путем активации сигнала «ОВ КНН», что характеризуется активным состоянием сигнала «КНН: Оперативный вывод» на выходе функции. 
\item
Алгоритм работы функции выполнен в соответствии с рисунком \ref{knn:img1}. В таблице \ref{knn:tbl1} приведены параметры, необходимые для настройки функции.

\vspace{3mm}
\begin{figure}[H]
\centering
\includegraphics[width=0.8\textwidth,height=0.8\textheight,keepaspectratio]{\fbpath/210. КНН ТН/_latex/img/knn.pdf}
\captionof{figure}{Функциональная схема логики КНН}\label{knn:img1}
\end{figure}

\end{enumerate}
